% Generated draft: requires mathematical and source review.
\documentclass[11pt]{article}
\usepackage[margin=25mm]{geometry}
\usepackage{fontspec}
\setmainfont{Latin Modern Roman}
\usepackage{amsmath,amssymb,mathtools,unicode-math}
\setmathfont{Latin Modern Math}
\usepackage{enumitem,needspace,xurl,hyperref}
\setlength{\parindent}{0pt}
\setlength{\parskip}{5pt}
\setlength{\emergencystretch}{3em}
\newcommand{\N}{\mathbb{N}}
\newcommand{\Z}{\mathbb{Z}}
\newcommand{\Q}{\mathbb{Q}}
\newcommand{\R}{\mathbb{R}}
\newcommand{\C}{\mathbb{C}}
\newcommand{\F}{\mathbb{F}}
\newcommand{\A}{\mathbb{A}}
\newcommand{\PP}{\mathbb{P}}
\newcommand{\E}{\mathbb{E}}
\newcommand{\eps}{\varepsilon}
\newcommand{\dd}{\,\mathrm{d}}
\begin{document}
\section*{UnsolvedMath ID: 30005149}
\textbf{Source problem code:} catalogue entry
\subsection{Definitions and mathematical statement}
The supplied record asks an expository and historical question, rather than specifying a conjecture with a definite mathematical conclusion. The following standard definitions provide context; their conventions are not additional hypotheses supplied by the source.

Let \(\mathbb K\in\{\R,\C\}\). A chain complex consists of \(\mathbb K\)-vector spaces \(V_k\) and linear maps \(\partial_k:V_k\to V_{k-1}\) satisfying
\[
\partial_{k-1}\circ\partial_k=0.
\]
Its cycles and boundaries are respectively
\[
Z_k=\ker\partial_k,\qquad B_k=\operatorname{im}\partial_{k+1}.
\]
Thus \(B_k\subseteq Z_k\), and its homology is \(H_k=Z_k/B_k\). The supplied excerpt instead says ``boundaries modulo cycles.'' That wording reverses the standard quotient and does not in general define a quotient space. The definition here explicitly uses cycles modulo boundaries.

A cochain complex has maps \(d^k:V^k\to V^{k+1}\) with \(d^{k+1}d^k=0\). In a standard analytic convention, a Hilbert complex consists of Hilbert spaces \(W^k\) over \(\mathbb K\) and densely defined closed linear operators
\[
d^k:D(d^k)\subseteq W^k\longrightarrow W^{k+1}
\]
satisfying
\[
d^k(D(d^k))\subseteq D(d^{k+1}),\qquad d^{k+1}d^k u=0\quad\text{for every }u\in D(d^k).
\]
These conditions alone do not impose closed ranges. The excerpt does not specify an analytic convention for Hilbert complexes or additional operator hypotheses. It mentions simplicial chain complexes associated with triangulations by simplices of a domain in \(\R^n\) as an illustrative example, without specifying a particular domain or triangulation.

The source question is: How did Hilbert complexes come to play a major role in numerical partial differential equations? In the supplied workshop context, this requests a conceptual and historical account of their relationship with numerical PDE. No particular PDE, spatial domain, boundary conditions, differential operators, discretization, error estimate, or quantified conclusion is specified. Consequently, the supplied evidence does not determine a precise theorem or conjecture to prove. Selecting such a target would require information beyond this record.
\subsection{Short English statement}
The record asks for an explanation of how Hilbert complexes became important in numerical PDE. It does not specify a theorem, PDE, boundary conditions, or numerical method. Standard complex definitions provide context, with an explicit correction of the excerpt's reversed description of homology. The intended mathematical target remains unspecified.
\subsection{Sources}
\begin{itemize}
\item Supplied local UnsolvedMath registry record, permanent ID \(30005149\), title ``Hilbert Complexes in Numerical Partial Differential Equations,'' exact source code \texttt{catalogue entry}, in the collection ``Oberwolfach Reports Open Problems.'' No individual proposer is supplied; the proposed year is \(2022\).
\item Supplied report citation: \emph{Hilbert Complexes: Analysis, Applications, and Discretizations} (2023), Oberwolfach Reports, DOI \texttt{10.4171/owr/2022/29}, pp.\ 1603--1659. \url{https://doi.org/10.4171/owr/2022/29}. The excerpt describes an introductory talk at the workshop held in June 2022 at Oberwolfach. The supplied citation year and workshop date are retained as given.
\item Supplied literature reference, ``OWR report'': \url{https://ems.press/content/serial-article-files/46966}.
\item Supplied literature reference, ``Arnold--Falk--Winther, finite element exterior calculus'': \url{https://doi.org/10.1137/1.9781611975543}. No further bibliographic details are supplied.
\item Literature assessment dated exactly \texttt{2026-08-22}: the supplied curation labels the record ``open,'' describes it as an open-ended foundational question with no single resolved answer, and states that Hilbert-complex methods are well established in numerical PDE. This is supplied dated curation, not an independently verified literature assessment. The background explicitly calls for independent verification of time-sensitive status and cited claims.
\item The registry labels the statement \texttt{corrected\_verified} and reports that only formula delimiters, notation, grammar, and typography were repaired while preserving mathematical content and scope. These are supplied editorial assertions, not independent mathematical verification of this draft or of the source excerpt.
\end{itemize}
\end{document}
